% =============================================================================
% Chapter 5: Sprint 2 --- Real-Time Collaboration and AI Core
% Assigned to: Asma Loudhaief
% =============================================================================

\chapter{Sprint 2 --- Real-Time Collaboration and AI Core}
\label{chap:sprint2}

% TODO: Assigned to Asma Loudhaief
% This chapter covers Sprint 2 (21 User Stories), focusing on low-latency WebSockets,
% live cursor tracking, the decoupled ai-service, multi-model LLM workflows,
% rate limiting, and Resilience4j circuit breaking.

\section{Introduction}
\label{sec:ch5_intro}
Sprint 2 introduced the two defining technological capabilities of MindSpark: bi-directional multi-user real-time interaction and an integrated Generative AI assistance engine.

\section{Sprint Goal \& Backlog}
\label{sec:ch5_goal_backlog}
\subsection{Sprint Goal Definition}
\label{subsec:ch5_sprint_goal}
% Formal statement of the Sprint 2 Goal centered on real-time synchronization and AI integration.

\subsection{Sprint 2 Backlog}
\label{subsec:ch5_sprint_backlog}
% Table of the 21 user stories selected for Sprint 2 with story points and assigned developers.

\section{Sprint Analysis}
\label{sec:ch5_analysis}
\subsection{Sprint Use Case Diagram}
\label{subsec:ch5_use_case_diag}
% UML Use Case diagram for real-time collaboration and AI operations.
\placeholderfig{fig5_1_sprint2_use_case.png}{Sprint 2 Use Case Diagram}{fig:sprint2-use-case}

\subsection{Textual Use Case Descriptions}
\label{subsec:ch5_use_case_specs}
% Operational specifications for core scenarios:
% - UC05: Broadcast Live Cursors and Online Presence
% - UC06: Synchronize Sticky Note Drag and Object Locking
% - UC07: Free-Hand Canvas Drawing Broadcast
% - UC08: Execute AI Idea Generation & Expansion
% - UC09: Execute AI Semantic Clustering & Grouping
% - UC10: Conduct Structured Voting Sessions

\section{Sprint Design}
\label{sec:ch5_design}
\subsection{Sprint Class Diagram}
\label{subsec:ch5_class_diag}
% Updated class diagram showing Vote, VotePool, Drawing, and AI Service entity models.
\placeholderfig{fig5_2_sprint2_class_diagram.png}{Sprint 2 Domain Class Diagram}{fig:sprint2-class-diagram}

\subsection{Design Sequence Diagrams}
\label{subsec:ch5_sequence_diags}
% Sequence diagrams detailing:
% 1. WebSocket bi-directional handshake and message broadcast
% 2. AI request flow: Gateway -> AI-Service -> Gemini/Groq -> Circuit Breaker fallback
\placeholderfig{fig5_3_seq_ai_generation.png}{Sequence Diagram: AI Generation and Circuit Breaker Flow}{fig:seq-ai-generation}

\subsection{Activity and State-Transition Diagrams}
\label{subsec:ch5_activity_state}
% State-transition diagram for Resilience4j Circuit Breaker (Closed, Open, Half-Open)
% and activity diagram for voting session execution.

\section{Implementation \& Screenshots}
\label{sec:ch5_implementation}
\subsection{Technical Realization}
\label{subsec:ch5_tech_realization}
% Native Spring WebSocket handler, Angular RealtimeService with Signals,
% ai-service dual-provider REST clients (Gemini, Groq), and concurrency rate limiter.

\subsection{User Interface Demonstration}
\label{subsec:ch5_ui_screenshots}
% Screenshots demonstrating:
% - Real-time multi-cursor collaboration with participant name tags
% - AI Assistant panel executing semantic idea grouping
% - Voting heatmap on sticky notes
\placeholderfig{fig5_4_realtime_ai_screenshot.png}{Real-Time Multi-User Collaboration and AI Assistant Panel}{fig:sprint2-ui-screenshot}

\section{Sprint Review \& Retrospective}
\label{sec:ch5_review_retro}
\subsection{Sprint Review Outcomes}
\label{subsec:ch5_review}
% Review meeting results demonstrating live WebSocket collaboration and AI workflows.

\subsection{Sprint Retrospective Takeaways}
\label{subsec:ch5_retro}
% Reflection on managing slow upstream LLM responses, concurrency tuning, and team action items.

\section{Conclusion}
\label{sec:ch5_conclusion}
% Summary of Sprint 2 achievements and transition to Sprint 3.
